\section{La frontera entre las empresas de aplicaciones y las empresas de modelos se volverá difusa}

La sección anterior trató la pila del sistema de generación; esta sección analiza las fronteras entre empresas. Wang Guan (王冠) considera que la frontera entre las empresas de aplicaciones y las empresas de modelos se volverá difusa. La razón es que las empresas de modelos subirán hacia el producto, mientras que las empresas de aplicaciones bajarán hacia los datos, el DSL, el context y la optimización de modelos locales. La diferencia real no está en la etiqueta de la empresa, sino en quién posee el volante de datos y los escenarios de usuario más decisivos.

\begin{figure}[H]
\centering
\includegraphics[width=0.94\textwidth]{app-model-boundary.png}
\caption{Frontera entre las empresas de aplicaciones y las empresas de modelos: el modelo base y los datos generados por el propio producto forman volantes distintos. Diagrama conceptual propio, elaborado a partir de la conversación de 01:41:59--02:01:44.}
\end{figure}

\begin{knowledgebox}{Lectura de la figura: una frontera difusa no es una frontera que desaparece}
Las empresas de modelos poseen el modelo base, la capacidad de cómputo, las capacidades generales y la distribución por API; las empresas de aplicaciones poseen los escenarios, los datos generados por el propio producto, los flujos de trabajo y el valor para el usuario. Lo decisivo en la competencia futura será quién logre conectar las capacidades del modelo con los datos del escenario en un ciclo cerrado.
\end{knowledgebox}

\subsection{Por qué las empresas de modelos suben hacia el producto}

Esta subsección examina primero por qué las empresas de modelos se desplazan hacia la capa de aplicaciones. Si las capacidades del modelo se quedan solo en la API, quedan demasiado lejos de la retroalimentación de los usuarios y del techo de ingresos; un punto de entrada de producto puede convertir las capacidades del modelo en una relación comercial de mayor valor.

Si una empresa de modelos solo vende API, queda fácilmente restringida por el precio y el costo de cómputo; subir hacia el producto le permite obtener el punto de entrada de los usuarios, retroalimentación real e ingresos más altos. ChatGPT, Claude Code, Gemini y otros demuestran que las empresas de modelos entran directamente en la capa de aplicaciones. Sobre todo en escenarios como Coding, las empresas de modelos ya controlan de origen los datos y el punto de entrada de los desarrolladores, por lo que les resulta más fácil hacerlo por su cuenta.

\subsection{Por qué las empresas de aplicaciones deben bajar hacia los modelos y los datos}

A continuación se examina el mismo problema desde la perspectiva de las empresas de aplicaciones. Si las empresas de modelos suben hacia el producto, las empresas de aplicaciones no pueden quedarse en la capa superficial, sino que deben construir hacia abajo datos, métodos y flujos de trabajo.

Si una empresa de aplicaciones solo hace la UI, también será copiada por las empresas de modelos. Debe construir hacia abajo el DSL, el context, los datos del dominio, la evaluación, la retroalimentación y la optimización local. No necesariamente tiene que entrenar un gran modelo general, pero sí necesita que su sistema de generación cuente con datos generados por el propio producto que otros no puedan obtener.

\begin{importantbox}{La ventaja defensiva de las empresas de aplicaciones}
La ventaja defensiva de una empresa de aplicaciones no es solo la interfaz, sino «el volante de datos dentro del escenario»: la intención del usuario, el proceso de producción, los ejemplos fallidos, las recipe, la retroalimentación comercial y la evaluación continua.
\end{importantbox}

\subsection{Resumen de la sección}

La frontera entre las empresas de aplicaciones y las empresas de modelos se volverá difusa, pero las diferencias no desaparecerán. Las empresas de modelos parten de las capacidades; las empresas de aplicaciones, de los escenarios. Quien logre formar un volante de datos más sólido obtendrá una posición más estable en la era de los sistemas de generación.
